% 付録：階差型・総和型から微積分と近似計算へ

\section{一歩の変化を細かくすると,微積分になるか}\label{節：差分から始まる離散的な微積分}
  \BookFigDifference


前節では,本編で等比数列を作った置き換えを,行列の作用として捉え直した.
ここでは,\sectionref{節：階差型}の「差から元の数列を求める」という考え方を出発点にしよう.
その逆向きの操作は,\sectionref{節：総和型}で使った「和から差を取って一般項を取り出す」という操作である.
この二つを一緒に見ると,差分と和が微分と積分にどう対応するかが分かる.

以下では,最初の項を$a_0$と書くことがある.
数列$a_0,a_1,a_2,\ldots$の隣り合う項の差
\[
  \Delta a_n=a_{n+1}-a_n
\]
を一階差分という.
本編で使った$a_{n+1}-a_n$に,\,$\Delta$という記号と「差分」という名前を付けたのである.
たとえば$a_n=n^2$なら
\[
  \Delta a_n=2n+1,\qquad \Delta^2a_n=2.
\]
二次式から差を取ると一次式になり,もう一回差を取ると定数になる.
ここで$\Delta^2a_n$は,\,$\Delta a_{n+1}-\Delta a_n$の意味である.

本編の特性方程式型の別解でも,添字を一つずらして引くことで定数項$q$を消し,
\[
  \Delta a_{n+1}=p\,\Delta a_n
\]
という等比型を作った.
その計算も,差分によって扱いやすい量を取り出した例だった.

逆に,差分を足し合わせれば元の数列を求め直せる.
\[
  \begin{aligned}
    a_n-a_0
      &=(a_1-a_0)+\cdots+(a_n-a_{n-1})\\
      &=\sum_{k=0}^{n-1}\Delta a_k.
  \end{aligned}
\]
途中の項が打ち消され,最初と最後だけが残る.
本編の階差型で和を取った操作は,離散的な積分と見直せる.
総和型で$S_{n+1}-S_n=a_{n+1}$とした操作は,その逆である.

この見方は,\sectionref{節：階差等比複合型}で引く多項式$g(n)$を探したことにもつながる.
特に$p=1$なら,\,$g(n+1)-g(n)=f(n)$を満たす式を探す問題になる.
たとえば$S_n=\sum_{k=1}^n k^2$については
\[
  S_n-S_{n-1}=n^2,\qquad S_0=0.
\]
二次式の差を逆向きに戻すので,三次式
$S_n=An^3+Bn^2+Cn+D$を候補にする.
その差は
\[
  S_n-S_{n-1}
  =3An^2+(-3A+2B)n+(A-B+C).
\]
係数比較から$\displaystyle A=\frac{1}{3}$, $\displaystyle B=\frac{1}{2}$, $\displaystyle C=\frac{1}{6}$となり,\,$S_0=0$から$D=0$である.
この候補が同じ初期値と差の関係を満たすことにより
\[
  S_n=\frac{n(n+1)(2n+1)}6
\]
が確かめられる.
和の公式を覚えて使うだけでなく,「差を取ると目的の式に戻るものは何か」と考えて作ることができた.

ここからは微分と一次近似を使い,差分と連続的な変化の関係を見よう.
時刻$t_n=nh$の値を$a_n$とすると,一単位時間あたりの変化は
\[
  \frac{a_{n+1}-a_n}{h}
\]
である.
ある微分可能な関数$y$の値$a_n=y(t_n)$を記録したものなら,刻み幅$h$を小さくすると,この差分商は微分係数$y'(t_n)$に近づく.

逆に,微分方程式$y'=F(t,y)$を数値で解くには,今の点での傾きから
\[
  a_{n+1}=a_n+hF(t_n,a_n)
\]
と次の値を予測できる.
これをオイラー法という.
正確な連続解の値と一致するとは限らないが,接線の方向へ一歩ずつ進む漸化式で,曲線を近似している.

本編の等比型と比較しやすい例として,\,$y'=-y$, $y(0)=1$を考えよう.
連続解は$y(t)=e^{-t}$で,時間が進むと$0$に近づく.
オイラー法では
\[
  a_{n+1}=(1-h)a_n,\qquad a_0=1
\]
だから,\,$a_n=(1-h)^n$である.
本編の等比型で調べたように,\,$0$への収束条件は$|1-h|<1$, すなわち$0<h<2$になる.
$0<h<1$では正のまま減少するが,\,$1<h<2$では符号を交互に変えながら近づく.
さらに$h>2$では絶対値が増大する.
本当の解が減衰していても,刻み幅を大きくすると,近似計算の列はまったく違う動きになりうる.
本編で等比数列の公比を調べた知識は,数値計算が安定するかを判断する道具にもなる.

また,正の時刻$t$を$N$等分して$\displaystyle h=\frac{t}{N}$とすれば,時刻$t$における近似値は
\[
  a_N=\left(1-\frac{t}{N}\right)^N
  \longrightarrow e^{-t}.
\]
刻み幅を細かくすることと,近似値が連続解へ近づくことが,この例では基本的な指数関数の極限で確認できる.
一般の微分方程式でも同様の収束を期待するが,その保証には関数や刻み幅の条件を調べる必要がある.

次に,\sectionref{節：ニュートン法型}で扱った近似計算へ戻ろう.
本編では,接線と横軸の交点を次の値に選ぶことで
\[
  a_{n+1}=\frac12\left(a_n+\frac2{a_n}\right)
\]
を作り,\,$\sqrt2$への誤差がほぼ二乗されることを見た.
ここではその具体例から,なぜ一般のニュートン法でも誤差が二乗程度になるのかへ進む.

方程式$g(x)=0$に対するニュートン法の更新は
\[
  x_{n+1}=x_n-\frac{g(x_n)}{g'(x_n)}
\]
である.
$g(\alpha)=0$, $g'(\alpha)\ne0$とし,根$\alpha$の近くで$g$は二回連続微分可能とする.
本編のずれの置き換えと同じく,\,$e_n=x_n-\alpha$とおくと
\[
  e_{n+1}
  =\frac{g'(x_n)e_n-g(x_n)}{g'(x_n)}.
\]
分子は,接線による変化と実際の変化との差である.
二階微分$g''$の絶対値がこの近くで$M$以下なら
\[
  \begin{aligned}
    g'(x_n)e_n-g(x_n)
      &=\int_\alpha^{x_n}\{g'(x_n)-g'(t)\}\,dt,\\
    |g'(x_n)e_n-g(x_n)|
      &\le \frac{M}{2}|e_n|^2.
  \end{aligned}
\]
二つ目の式は,平均値の定理から
$|g'(x_n)-g'(t)|\le M|x_n-t|$として積分すれば得られる.
積分の向きが逆になる$x_n<\alpha$でも,絶対値を取れば同じ評価になる.
また,\,$g'(\alpha)\ne0$と連続性により,十分近くでは$|g'(x_n)|\ge m$となる正の定数$m$を選べる.
したがって
\[
  |e_{n+1}|\le\frac{M}{2m}|e_n|^2.
\]
誤差が十分小さければ,次の誤差はその二乗の定数倍以下になる.
いったん十分近くから始めれば近くにとどまり,急速に根へ近づくこともこの評価から分かる.
本編の$\sqrt2$の例で得た正確な式
\[
  e_{n+1}=\frac{e_n^2}{2a_n}
\]
は,この一般的な性質が特に見えやすい形で現れたものだった.

前節の不動点の見方でも確かめられる.
更新関数$\displaystyle N(x)=x-\frac{g(x)}{g'(x)}$について
\[
  N'(\alpha)
  =\left.\frac{g(x)g''(x)}{g'(x)^2}\right|_{x=\alpha}=0.
\]
つまり,根のところでは一次のずれの倍率が零になり,二次のずれが残る.
本編で等比型へ帰着させたときは一定倍率でずれを追ったが,ニュートン法ではずれが小さいほど,次に残る割合も小さくなるのである.

接線の傾きを毎回微分で求めるのが難しいなら,過去二点を結ぶ割線の傾きで代用できる.
その更新は
\[
  x_{n+1}
  =x_n-g(x_n)\frac{x_n-x_{n-1}}{g(x_n)-g(x_{n-1})}
\]
となり,割線法と呼ばれる.
分母が零でない範囲で用いる.
直前の二項を記録して次を作るという点は,前節の二成分の状態と共通している.
本編で与えられた漸化式を解いた経験は,このように近似計算の規則を作り,その収束を調べることにもつながる.
